% !PRINTED 409
\weberpage{409}
Für den Fall eines \gesperrt{geraden} $\Delta$ zerlegt man $k$ in

\[
v, \quad -\frac{\Delta}{2} + v, \qquad 0 < v < \frac{\Delta}{2}
\]

und wendet die Formel an:

\[
\left(\Delta,\ -\frac{\Delta}{2} + v\right) = -(\Delta,v),
\]

die sich aus den Sätzen des § 85 leicht ergibt.

Dann findet man zunächst:

\[
h = \frac{1}{2}\overset{v}{\Sigma}(\Delta,v),
\]

und wenn man die $v$ nochmals zerlegt in

\weq{12}{\begin{aligned}
&\mu, \quad -\frac{\Delta}{2} - \mu: \qquad 0 < \mu \leqq \frac{-\Delta}{4},\\
&h = \Sigma(\Delta,\mu).
\end{aligned}}

Die Formel für positive Diskriminanten läßt sich mit Hilfe
der Kreisteilungstheorie umformen, worauf wir hier nicht eingehen
wollen.

\begin{center}
§ 113.\quad \gesperrt{Die Anzahl der Geschlechter.}
\end{center}

Die Formel (5) des vorigen Paragraphen gestattet die Bestimmung
 der genauen Anzahl der Geschlechter nach \gesperrt{Dirichlet}.
Wir verstehen unter $\delta$ irgend einen Stammteiler von $D$ und setzen:

\weq{1}{\begin{aligned}
&F(z) = \frac{(\delta,z)}{z^s}, \quad \text{wo } z \text{ relativ prim zu } D \text{ ist,}\\
&F(z) = 0, \qquad \text{wenn } z \text{ und } D \text{ einen gemeinsamen Teiler haben.}
\end{aligned}}

Dann ist $(\delta,N\mathfrak{a})$ einer der Charaktere des durch $\mathfrak{a}$ bestimmten
Geschlechtes. Bezeichnen wir diesen Charakter also mit $(\delta,A)$,
so folgt aus § 112, (5):

\weq{2}{\overset{A}{\Sigma}(\delta,A)\overset{\mathfrak{a}}{\Sigma}\frac{1}{(N\mathfrak{a})^s} = \Sigma\frac{(\delta,m)}{m^s}\Sigma\frac{(\delta',n)}{n^s}.}

Darin ist $\delta' = \delta\Delta$ und kann, da $n$ relativ prim zu $D$ ist,
auch durch den \gesperrt{Stamm} von $\delta\Delta$ ersetzt werden (nach der symbolischen
 Multiplikation in § 104). Ist $\delta = 1$, so ist $\delta' = \Delta$,
und ist $\delta = \Delta$, so ist $\delta' = 1$.

Von den beiden Fällen $\delta = 1$, $\delta = \Delta$ abgesehen, behalten
also die Summen auf der rechten Seite von (2) für $s = 1$ einen
endlichen Wert (die nach dem vorigen Paragraphen durch Klassenzahlen
 der Diskriminanten $\delta$, $\delta'$ ausdrückbar sind).

Auf der linken Seite von (2) durchläuft in der Summe nach $\mathfrak{a}$
das Zeichen $\mathfrak{a}$ die Gesamtheit der Ideale der einzelnen \gesperrt{Klassen},
%CONT>>
